\documentclass[11pt,a4paper]{article}
\usepackage[utf8]{inputenc}
\usepackage[T1]{fontenc}
\usepackage[spanish]{babel}
\usepackage[margin=2.5cm]{geometry}
\usepackage{amsmath,amssymb}
\usepackage{color}
\usepackage[hidelinks]{hyperref}
\setlength{\parskip}{5pt}
\setlength{\parindent}{0pt}
\title{Computadora de vuelo del cohete de agua\\\large Estimación de la velocidad vertical $V_r$}
\author{}
\date{}
\begin{document}
\hypersetup{pageanchor=false}
\begin{titlepage}
\vspace*{3cm}
\begin{center}
{\Large Computadora de vuelo del cohete de agua\par}
\vspace{2cm}
{\color{red}\Huge\bfseries CONFIDENCIAL\par}
\vspace{0.3em}
{\color{red}\Large\bfseries Acceso exclusivo para\par
Itzel, Freddy y Salvador\par}
\end{center}
\vfill
\end{titlepage}
\hypersetup{pageanchor=true}

\maketitle
\vspace{-1.5em}
La velocidad vertical indica qué tan rápido cambia la altura del cohete. Tomaremos el sentido ascendente como positivo: una velocidad positiva corresponde al ascenso y una negativa al descenso. En este documento llamaremos \textbf{$V_r$ (velocidad vertical Real)} al resultado de combinar el barómetro y el acelerómetro. El nombre identifica la variable del proyecto; su valor es una estimación de la velocidad física, no una medición exacta.

\section{Velocidad vertical a partir del barómetro}
El barómetro mide presión atmosférica. Como la presión disminuye con la altura, sus variaciones permiten estimar el desplazamiento vertical. A partir del equilibrio hidrostático y del modelo de aire ideal,
\begin{equation}
 \frac{dp}{dh}=-\rho g,\qquad \rho=\frac{p}{R_a T},
\end{equation}
se obtiene, suponiendo temperatura media y gravedad aproximadamente constantes,
\begin{equation}
 h_b(t)=\frac{R_a\bar T}{g}\ln\left(\frac{p_0}{p(t)}\right).
\end{equation}
Aquí $h_b$ es la altura barométrica respecto al punto de lanzamiento, $p_0$ la presión de referencia medida antes de despegar, $p(t)$ la presión actual, $\bar T$ la temperatura media del aire en kelvin y $R_a\approx287\,\mathrm{J/(kg\,K)}$. Se toma $g\approx9.81\,\mathrm{m/s^2}$. Las presiones deben expresarse en las mismas unidades.

La velocidad es la derivada de la altura. Para dos mediciones consecutivas separadas por $\Delta t_k=t_k-t_{k-1}$,
\begin{equation}
 V_{b,k}=\frac{h_{b,k}-h_{b,k-1}}{\Delta t_k}
 =\frac{R_a\bar T}{g\Delta t_k}\ln\left(\frac{p_{k-1}}{p_k}\right).
\end{equation}
Si la presión disminuye, $V_{b,k}$ es positiva. Esta diferencia aproxima la velocidad media del intervalo; por sí sola no entrega la velocidad instantánea exacta en $t_k$.

Derivar también amplifica el ruido: si dos alturas tienen errores independientes con desviación estándar $\sigma_h$, entonces
\begin{equation}
 \sigma_{V_b}\approx\frac{\sqrt{2}\,\sigma_h}{\Delta t_k}.
\end{equation}
Por ello se necesita suavizar las mediciones o estimar la pendiente sobre varias muestras, aceptando cierto retardo. El barómetro debe medir presión estática exterior: las perturbaciones del flujo y la ventilación del alojamiento pueden introducir errores que el promedio no elimina.

\newpage
\section{Velocidad vertical a partir del acelerómetro}
El acelerómetro mide fuerza específica: su lectura incluye el efecto de sostener al sensor contra la gravedad. En reposo su magnitud es aproximadamente $g$; en caída libre ideal es aproximadamente cero. Para obtener la aceleración vertical de trayectoria hay que corregir sesgos, orientar la lectura y compensar la gravedad:
\begin{equation}
 a_{z,k}=\mathbf e_z^{\mathsf T}\mathbf R_{NB,k}
 (\mathbf f_{m,k}-\mathbf b_f)-g.
\end{equation}
$\mathbf f_m$ es la lectura en $\mathrm{m/s^2}$, previamente corregida en escala; $\mathbf b_f$ es su sesgo; $\mathbf R_{NB}$ transforma del marco del sensor al marco local, y $\mathbf e_z=[0\;0\;1]^{\mathsf T}$ selecciona la vertical ascendente.

Si el eje $z$ del sensor permanece vertical y apunta hacia arriba, la expresión se reduce a $a_z=f_z-b_z-g$. Si el cohete se inclina, se necesita conocer su orientación, por ejemplo mediante el giróscopo. Durante el vuelo dinámico, el acelerómetro por sí solo no permite separar de forma general inclinación y aceleración de traslación.

La integración permite actualizar la velocidad. Usando la aceleración media de dos muestras,
\begin{equation}
 V_{a,k}=V_{r,k-1}+\frac{a_{z,k-1}+a_{z,k}}{2}\Delta t_k.
\end{equation}
$V_{a,k}$ es la predicción inercial del instante actual, obtenida desde la estimación combinada anterior. Antes del despegue, con el cohete en reposo, se establece $V_{r,0}=0$. La integración responde rápidamente, pero un sesgo residual constante $b_a$ acumula un error $\delta V=b_a t$.

\section{Combinación para obtener $V_r$}
Se propone comenzar con un filtro complementario: el acelerómetro describe los cambios rápidos y el barómetro corrige gradualmente la deriva de la velocidad integrada. Cuando existe una nueva medición barométrica válida,
\begin{equation}
 \boxed{V_{r,k}=\alpha_k V_{a,k}+(1-\alpha_k)V_{b,k}},
 \qquad \alpha_k=\frac{\tau}{\tau+\Delta t_k}.
\end{equation}
La constante $\tau>0$, expresada en segundos, determina la rapidez de la corrección. Un valor grande confía más en la predicción inercial; uno pequeño introduce más corrección barométrica y también más ruido. Su valor se elegirá a partir de mediciones, sin fijarlo todavía.

Estas ecuaciones suponen muestras sincronizadas y retardos pequeños respecto a los cambios del vuelo. Si el barómetro se actualiza más lentamente, se propaga la velocidad con el acelerómetro entre lecturas y se corrige solo al recibir una lectura nueva; el peso de corrección debe corresponder al intervalo entre correcciones. Una muestra barométrica repetida no constituye información nueva.

Así, $V_{r,k}$ representa la mejor estimación disponible de la velocidad vertical en $t_k$ con este método. Su precisión depende de la calibración, la orientación, el ruido y el retardo de los sensores; deberá comprobarse experimentalmente. Combinar ambas fuentes reduce sus limitaciones complementarias, pero no garantiza una velocidad exacta.
\newpage
\section{Ejemplo: ascenso con inclinación lateral}
Consideremos dos instantes separados por $\Delta t=0.10\,\mathrm{s}$ durante el ascenso. El cohete realiza una rotación suave y controlada hacia el lateral $+x$ del marco local. La inclinación de su eje longitudinal pasa de $10^\circ$ a $20^\circ$ respecto a la vertical, mediante una rotación alrededor del eje $y$. El carácter controlado de esta maniobra es una hipótesis del ejemplo, no una función de control definida aquí.

Los valores son ilustrativos. Suponemos lecturas sincronizadas y calibradas, orientación conocida por una estimación independiente apoyada en el giróscopo, y un acelerómetro próximo al centro de masa, de modo que se desprecia el efecto de su desplazamiento respecto a ese punto. Las componentes indicadas son fuerza específica en los ejes del sensor, expresada en $\mathrm{m/s^2}$.
\begin{center}
\begin{tabular}{c|c|c|c}
Instante & Presión & Inclinación & Lectura $(f_x,f_y,f_z)$ \\
\hline
$t_1$ & $100000\,\mathrm{Pa}$ & $10^\circ$ & $(1.00,\ 0.00,\ 15.00)$ \\
$t_2=t_1+0.10\,\mathrm{s}$ & $99990.5\,\mathrm{Pa}$ & $20^\circ$ & $(2.00,\ 0.00,\ 14.00)$
\end{tabular}
\end{center}
El movimiento se limita al plano $xz$ para facilitar el cálculo; por eso $f_y=0$. Tomamos $\bar T=293.15\,\mathrm{K}$ y una velocidad previamente estimada $V_r(t_1)=8.00\,\mathrm{m/s}$. Esta condición inicial es necesaria: dos lecturas de aceleración solo proporcionan el cambio de velocidad.

\subsection{Aporte del barómetro}
Con las dos presiones, la velocidad barométrica resulta
\begin{equation}
 V_b=\frac{287(293.15)}{9.81(0.10)}
 \ln\left(\frac{100000}{99990.5}\right)
 \approx8.148\,\mathrm{m/s}.
\end{equation}
La caída de presión equivale a un incremento de altura aproximado de $0.815\,\mathrm{m}$. Esta velocidad corresponde al promedio del intervalo y es positiva porque el cohete asciende.

\subsection{Proyección de las componentes del acelerómetro}
Para la rotación lateral definida, la transformación del marco del sensor al marco local es
\begin{equation}
 \mathbf R_{NB}(\theta)=
 \begin{bmatrix}
 \cos\theta&0&\sin\theta\\
 0&1&0\\
 -\sin\theta&0&\cos\theta
 \end{bmatrix}.
\end{equation}
Por tanto, las aceleraciones de trayectoria lateral y vertical son
\begin{equation}
 a_x=f_x\cos\theta+f_z\sin\theta,\qquad
 a_z=-f_x\sin\theta+f_z\cos\theta-g.
\end{equation}
La componente $f_z$ del sensor no es la componente vertical local cuando el cuerpo está inclinado. Además, la inclinación del cuerpo y la dirección de la aceleración no tienen por qué coincidir.

\newpage
En el primer instante se obtiene
\begin{align}
 a_{z,1}&=-1.00\sin10^\circ+15.00\cos10^\circ-9.81
 \approx4.788\,\mathrm{m/s^2},\\
 a_{x,1}&=1.00\cos10^\circ+15.00\sin10^\circ
 \approx3.590\,\mathrm{m/s^2}.
\end{align}
En el segundo instante,
\begin{align}
 a_{z,2}&=-2.00\sin20^\circ+14.00\cos20^\circ-9.81
 \approx2.662\,\mathrm{m/s^2},\\
 a_{x,2}&=2.00\cos20^\circ+14.00\sin20^\circ
 \approx6.668\,\mathrm{m/s^2}.
\end{align}
Los vectores de aceleración en el marco local son, aproximadamente,
\begin{equation}
 \mathbf a_{N,1}=\begin{bmatrix}3.590\\0\\4.788\end{bmatrix}\,\mathrm{m/s^2},
 \qquad
 \mathbf a_{N,2}=\begin{bmatrix}6.668\\0\\2.662\end{bmatrix}\,\mathrm{m/s^2}.
\end{equation}
Existe aceleración lateral, pero solo $a_z$ interviene en la actualización de velocidad vertical. Si se restara directamente $g$ a $f_{z,2}$, se obtendría $4.19\,\mathrm{m/s^2}$ en lugar de $2.662\,\mathrm{m/s^2}$: se estaría ignorando la orientación.

\subsection{Predicción inercial y velocidad combinada}
Integrando las componentes verticales con la regla del trapecio,
\begin{equation}
 V_a(t_2)=8.00+\frac{4.788+2.662}{2}(0.10)
 \approx8.373\,\mathrm{m/s}.
\end{equation}
Para ilustrar la combinación elegimos $\tau=0.40\,\mathrm{s}$, de modo que
\begin{equation}
 \alpha=\frac{0.40}{0.40+0.10}=0.80.
\end{equation}
Esta elección es didáctica; no representa un ajuste validado. Con los valores calculados antes de redondear,
\begin{equation}
 V_r(t_2)=0.80\,V_a(t_2)+0.20\,V_b
 \approx\boxed{8.328\,\mathrm{m/s}}.
\end{equation}
El resultado indica que el cohete sigue ascendiendo a una velocidad vertical estimada de aproximadamente $8.33\,\mathrm{m/s}$, aunque su aceleración también tiene una componente lateral. No es la magnitud de la velocidad total del cohete.

En este cálculo, la predicción inercial corresponde a $t_2$, mientras que el barómetro aporta un promedio entre $t_1$ y $t_2$. La combinación aproxima la velocidad actual bajo la hipótesis de un intervalo corto; esa diferencia temporal introduce error durante las aceleraciones. Los decimales mostrados permiten seguir la aritmética, pero no implican esa precisión física en los sensores.
\newpage
\section{Tiempo hasta velocidad vertical cero con resistencia del aire}
Durante el ascenso, la computadora utilizará la velocidad estimada $V_{r,k}$ y la aceleración vertical de trayectoria obtenida del acelerómetro. No se impondrá $a_z=-g$ como valor permanente: la medición permite incorporar el efecto actual del empuje y de las fuerzas aerodinámicas. La estimación requiere lecturas válidas, calibradas y corregidas por orientación.

\subsection{Qué mide el acelerómetro durante el vuelo}
El acelerómetro del MPU6050 mide fuerza específica, no la aceleración gravitatoria directamente. En caída libre ideal, sin resistencia del aire, su lectura es aproximadamente cero aunque el cohete acelere hacia abajo a $g$. En reposo sobre una superficie, la fuerza de soporte produce una lectura de magnitud cercana a $g$.

Definimos la componente vertical local de fuerza específica como
\begin{equation}
 f_{z,N,k}=\mathbf e_z^{\mathsf T}\mathbf R_{NB,k}
 (\mathbf f_{m,k}-\mathbf b_f),
 \qquad \boxed{a_{z,k}=f_{z,N,k}-g}.
\end{equation}
La lectura se supone corregida en escala y expresada en $\mathrm{m/s^2}$. Si el cohete está inclinado, $f_{z,N}$ combina sus tres ejes: no debe sustituirse por la componente $z$ cruda del sensor. Para la rotación lateral del ejemplo,
\begin{equation}
 a_z=-f_x\sin\theta+f_z\cos\theta-g.
\end{equation}
En un modelo de traslación referido al centro de masa,
\begin{equation}
 f_{z,N}=\frac{F_{T,z}+F_{\mathrm{aero},z}}{m},\qquad
 a_z=\frac{F_{T,z}+F_{\mathrm{aero},z}}{m}-g.
\end{equation}
$F_{T,z}$ es el empuje vertical y $F_{\mathrm{aero},z}$ la resultante aerodinámica vertical. Después del empuje, el sensor sigue registrando el efecto de la aerodinámica. Durante ascenso en aire quieto y con solo arrastre, ese efecto es negativo y $a_z$ es más negativa que $-g$.

\subsection{Cómo expresar el efecto del viento y el arrastre}
El arrastre depende de la velocidad respecto al aire. Si $\mathbf v_N$ es la velocidad del cohete respecto al suelo y $\mathbf w_N$ la velocidad del viento,
\begin{equation}
 \mathbf u=\mathbf v_N-\mathbf w_N,\qquad
 \mathbf F_D=-\tfrac12\rho C_D A\,\|\mathbf u\|\mathbf u.
\end{equation}
Aquí $\rho$ es la densidad del aire, $C_D$ el coeficiente de arrastre y $A$ su área de referencia. Es un modelo de arrastre opuesto al flujo relativo; una fuerza de sustentación o normal al flujo requeriría un término adicional.\footnote{NASA Glenn Research Center: \href{https://www1.grc.nasa.gov/beginners-guide-to-aeronautics/modern-drag-equation/}{Modern Drag Equation}.}
La componente vertical del modelo es
\begin{equation}
 a_z=\frac{F_{T,z}}{m}-g-
 \frac{\rho C_D A}{2m}\|\mathbf v_N-\mathbf w_N\|(V_z-w_z).
\end{equation}
El viento lateral modifica la rapidez relativa y, con ello, el arrastre; el viento vertical también puede cambiar el signo de su componente vertical. La sola estimación $V_r$ no determina el viento ni toda la velocidad relativa.

\newpage
\subsection{Predicción con la aceleración vertical medida}
Para actualizar $V_r$ se utilizará la aceleración vertical corregida durante el vuelo. Para predecir el apogeo, además se necesita estar ascendiendo y desacelerando. Sea $\bar a_{z,k}$ una estimación filtrada de esa aceleración, calculada con mediciones recientes; la barra indica suavizado temporal, no conocimiento de la aceleración futura.

Como aproximación local de movimiento rectilíneo uniformemente acelerado (MRUA), se supone que $\bar a_{z,k}$ permanece constante durante la predicción:
\begin{equation}
 V_z(t_k+\Delta t)\approx V_{r,k}+\bar a_{z,k}\Delta t.
\end{equation}
Al imponer velocidad vertical cero,
\begin{equation}
 \boxed{\Delta t_0\approx-\frac{V_{r,k}}{\bar a_{z,k}}},
 \qquad V_{r,k}>0,\quad \bar a_{z,k}<0.
\end{equation}
$\Delta t_0$ es el tiempo restante, mientras que $t_k+\Delta t_0$ es el instante estimado de apogeo. Si la aceleración es positiva, el modelo no predice una detención futura durante el ascenso. Si es cero o muy cercana a cero, no proporciona un tiempo finito fiable. En las lecturas del ejemplo anterior, $a_{z,2}=+2.662\,\mathrm{m/s^2}$, por lo que todavía no corresponde aplicar esta predicción de apogeo.

\subsection{Ejemplo después de terminar el empuje}
Supongamos ahora un instante distinto, de ascenso sin empuje, con $V_r=8.328\,\mathrm{m/s}$. Las lecturas del acelerómetro del MPU6050, calibradas, proyectadas a la vertical local y filtradas, dan $\bar f_{z,N}=-2.00\,\mathrm{m/s^2}$. Este valor proviene del sensor y todavía no tiene restada la gravedad. Al compensarla,
\begin{equation}
 \bar a_z=-2.00-9.81=-11.81\,\mathrm{m/s^2}.
\end{equation}
La estimación local del tiempo restante es
\begin{equation}
 \Delta t_0\approx\frac{8.328}{11.81}\approx0.705\,\mathrm{s}.
\end{equation}
Si se ignorara el arrastre se obtendría $8.328/9.81\approx0.849\,\mathrm{s}$. La diferencia refleja la desaceleración adicional que sí aparece en la medición. No debemos volver a restar un arrastre calculado a $\bar a_z$: hacerlo contaría dos veces la misma fuerza. El modelo aerodinámico sirve para interpretar o predecir la evolución, no para duplicar una contribución ya medida.

\subsection{Alcance de la estimación}
El acelerómetro observa el efecto actual del aire, pero no determina cómo cambiará hasta el apogeo. En ascenso puramente vertical, sin viento, sin empuje y con masa y parámetros aerodinámicos constantes, el modelo es
\begin{equation}
 \dot V_z=-g-\beta V_z^2,\qquad \beta=\frac{\rho C_D A}{2m}.
\end{equation}
Al disminuir la velocidad, disminuye el arrastre: mantener la desaceleración inicial constante tiende a subestimar el tiempo restante en este caso. Por eso $0.705\,\mathrm{s}$ es una extrapolación local, no un tiempo exacto. Se recalculará con cada actualización válida, considerando el retardo del filtrado.

Si el sensor está alejado del centro de masa, su lectura incluye aceleraciones tangenciales y centrípetas por la rotación; deben compensarse o demostrarse despreciables. Las lecturas saturadas o la orientación incierta invalidan la predicción. Finalmente, velocidad vertical cero no implica velocidad lateral cero, y la predicción temporal no sustituye la confirmación medida del apogeo.

\newpage
\section{Identificación de las lecturas del MPU6050}
Para evitar ambigüedades, distinguiremos el dato entregado por el acelerómetro del MPU6050 de la aceleración vertical que calcula la computadora de vuelo. Filtrar significa suavizar el ruido; no significa eliminar la gravedad ni corregir la inclinación.

\subsection{Del dato del sensor a la componente vertical}
El acelerómetro entrega tres mediciones digitales con signo, que denotamos $q_x,q_y,q_z$, en los ejes del sensor. Son lecturas de 16 bits; su sensibilidad depende del rango seleccionado.\footnote{TDK InvenSense, \href{https://invensense.tdk.com/wp-content/uploads/2015/02/MPU-6000-Register-Map1.pdf}{MPU-6000/MPU-6050 Register Map and Descriptions}, sección de mediciones del acelerómetro.}
Para una sensibilidad $S$ expresada en cuentas por $g$, la conversión es
\begin{equation}
 f_{m,i}=\frac{q_i}{S}\,g_0,\qquad i\in\{x,y,z\},
 \qquad g_0=9.80665\,\mathrm{m/s^2}.
\end{equation}
Por ejemplo, para rango $\pm2g$, $S=16384$ cuentas por $g$. Esta escala solo corresponde a ese rango. Si la lectura ya está en $\mathrm{m/s^2}$, no se convierte de nuevo. $g_0$ es la referencia de conversión; $g$ es la gravedad local, aproximada aquí por $9.81\,\mathrm{m/s^2}$.

Después se corrige el sesgo, se transforma a la vertical local y se filtra la señal proyectada:
\begin{align}
 f_{z,N,k}&=\mathbf e_z^{\mathsf T}\mathbf R_{NB,k}
 (\mathbf f_{m,k}-\mathbf b_f),\\
 \bar f_{z,N,k}&=\mathcal F\{f_{z,N}\}_k,\\
 \boxed{\bar a_{z,k}}&=\boxed{\bar f_{z,N,k}-g}.
\end{align}
$\mathcal F$ representa un filtro de suavizado que conserva el valor constante de la señal. La orientación debe corresponder al instante de cada lectura. $\bar f_{z,N}$ es fuerza específica vertical filtrada, mientras que $\bar a_z$ es aceleración vertical de trayectoria filtrada y compensada por gravedad. Es esta última la que se usa para predecir el tiempo restante.

\subsection{Qué significa el valor del ejemplo}
En el ejemplo de ascenso sin empuje se supuso
\begin{equation}
 \underbrace{\bar f_{z,N}=-2.00\,\mathrm{m/s^2}}_{\text{antes de restar gravedad}}
 \quad\Longrightarrow\quad
 \underbrace{\bar a_z=-11.81\,\mathrm{m/s^2}}_{\text{después de restar gravedad}}.
\end{equation}
El signo negativo es esencial: la fuerza específica vertical apunta hacia abajo. Ese $-2.00$ es un valor derivado de las lecturas del MPU6050, ya calibradas, proyectadas y filtradas; no es necesariamente la lectura de su eje $z$ ni una aceleración con gravedad ya compensada.

Si, en cambio, el dato disponible fuera $\bar a_z=-2.00\,\mathrm{m/s^2}$ \textbf{después de restar la gravedad}, se usaría directamente: no se volvería a restar $g$. Para $V_r=8.328\,\mathrm{m/s}$, la extrapolación sería $\Delta t_0=8.328/2.00=4.164\,\mathrm{s}$. Sería otro estado físico: no el ascenso sin empuje en aire quieto con solo arrastre del ejemplo anterior.

\end{document}

